\sisection{S7. Software, environment and the released audit tool}

\sisubsection{S7.1 Environment and report build}

The v5 report layer uses CPython 3.12.11 and the pinned \path{requirements.lock}. The lock file was
produced by \texttt{uv pip compile requirements.txt --python-version 3.12 --output-file
requirements.lock} and includes transitive dependencies. Table~\ref{tab:s7env} lists the
interpreter and the components imported directly by the analysis and report code.

\begin{table}[htbp]
\centering
\caption{Release interpreter and directly imported components. Package versions come from
\texttt{requirements.lock}; the interpreter is also pinned in \texttt{environment.yml}.}
\label{tab:s7env}
\small
\begin{tabular}{ll}
\toprule
Component & Version \\
\midrule
CPython & 3.12.11 \\
NumPy \cite{numpy2020} & 1.26.4 \\
pandas & 2.3.2 \\
SciPy \cite{scipy2020} & 1.16.2 \\
scikit-learn & 1.7.2 \\
RDKit \cite{rdkit2025} & 2025.03.6 \\
Matplotlib \cite{matplotlib2007} & 3.10.6 \\
Biopython & 1.87 \\
openpyxl & 3.1.5 \\
xlrd & 2.0.2 \\
pytest & 9.1.1 \\
\bottomrule
\end{tabular}
\end{table}

The normative local environment and the report build are:

\begin{verbatim}
uv venv --python 3.12.11 .venv
uv pip install --python .venv/bin/python -r requirements.lock
make PYTHON=.venv/bin/python all
\end{verbatim}

\path{Dockerfile} is the independent packaging check rather than a second dependency
specification. It starts from \texttt{python:3.12.11-slim}, adds a TeX Live subset, installs the
same \path{requirements.lock} through \texttt{pip}, and runs \texttt{make all}.

The default target is a report-layer build over frozen results. \texttt{make all} invokes
\texttt{submission-v5}: \texttt{verify-v5-numbers} regenerates and stages the submission figures,
including the graphical abstract, verifies the v5 report manifest, and runs the v5 consistency
and aggregate-reconstruction tests. \texttt{manuscript-v5} and
\texttt{supplement-v5} then compile the two PDFs twice with
\texttt{SOURCE\_DATE\_EPOCH=1786320000}. This build does not rerun every scientific producer,
create a release tag or mint a DOI. PDF bytes can still differ across TeX distributions even
when the text and page count agree.

The included \path{.github/workflows/v5-report.yml} is configured to run the complete
\texttt{release-check-v5} target in a clean Linux checkout and retain the tested PDFs and staged
artwork. This continuous-integration layer becomes meaningful only when every report input is
present in the committed release candidate.

The v5 manifest, \path{results/v5_report_source_manifest.csv}, records the exact local source,
frozen-result and build-contract files consumed by this report layer. It does not claim that the
upstream scientific producers were rerun. The strict public-core audit is separate. Its manifest,
\path{results/public_core_source_manifest.csv}, has 149 records: 111 inputs, 37 code files and one
output. \texttt{make verify-public-manifest} re-hashes the manifest and every registered file
without rewriting them. Some result bundles also carry \path{output_checksums.json}; coverage is
partial, so bundle summaries and metadata remain part of the frozen-output record.

Stochastic stages use explicit per-analysis seeds rather than one global value.
Table~\ref{tab:s7seeds} lists selected support and scorer seeds behind the reported results;
experimental-permutation seeds are specified with their endpoint contracts in S6 and in the
producer metadata. Changing a seed keeps the design but changes the replicate rows.

\begin{table}[htbp]
\centering
\caption{Selected support and scorer seeds for analyses reported in the main text. Sources are the metadata or
summary files in \texttt{calibration\_panel\_recovery}, \texttt{biological\_core\_modes} and
\texttt{nonvina\_scorer\_transport}; holdout seeds are also set in their producer script.}
\label{tab:s7seeds}
\small
\begin{tabular}{lr}
\toprule
Analysis & Seed \\
\midrule
Calibration recovery, DOCKSTRING-58 & 20260803 \\
Calibration recovery, Docking-44 & 20260804 \\
Chemical-group holdout, DOCKSTRING-58 & 20260805 \\
Chemical-group holdout, Docking-44 & 20260806 \\
DOCKSTRING 15{,}000-row working support & 71 \\
Five-mode core & 20260808 \\
Fixed-pose ligand subset & 20260802 \\
Fixed-pose cluster bootstrap, Bemis--Murcko & 20260805 \\
Fixed-pose cluster bootstrap, Butina & 20260806 \\
ODDT internal training & 1 \\
\bottomrule
\end{tabular}
\end{table}

Fixed-pose source rescoring used a separate environment. ODDT 0.7 supplies RF-Score, NNScore and
PLECscore \cite{oddt2015}; the released workflow runs its worker under CPython 3.9.23 from
conda-forge. That stage executes
\path{analysis/nonvina_scorer_oddt_worker.py} as a subprocess under the 3.9 interpreter, writes a
CSV of scores, and the spectral statistics are computed under the 3.12.11 environment from that
CSV. Both interpreter strings are written into the \texttt{toolchain} block of
\path{results/nonvina_scorer_transport/summary.json}. The 3.9 environment carries SciPy 1.13.1,
whose \texttt{scipy.sparse.vstack} no longer accepts a map object; the worker patches
\texttt{sparse\_vstack} in two ODDT namespaces to materialise the iterator into a list before
stacking, which the same record states has no numerical effect. The smina binary
\cite{smina2013} is identified by version string \texttt{smina 2020.12.10 conda-forge:b08c07c}
and by SHA-256, and it is not redistributed.

The report build therefore verifies and renders released results; it is not a source-level
clean-clone regeneration claim. \path{external_sources.csv} inventories the external datasets and
assay resources used by v5, with access, redistribution, licence, integrity and provenance fields.
The fixed-pose scorer toolchain additionally depends on ODDT training resources, DOCKSTRING
receptor PDBQTs and the smina executable identified in
\path{results/nonvina_scorer_transport/summary.json}; that summary is authoritative for their
versions and checksums. Core score matrices and several public assay tables are included. Other
inputs, including PDSP, SPD, KiRHub, PKIS1 and HotSpot source data, are not all redistributed. The
exact boundary for each analysis is in its result summary or README. The primary and complete-case PDSP paired-QAP tables and
support-threshold sensitivity can nevertheless be reconstructed from the released aggregate
target-pair table and Docking-44 matrix with
\path{analysis/reproduce_pdsp_derived_qap.py}; this does not reconstruct the upstream
compound-level PDSP preprocessing.

Fixed-pose rescoring additionally requires the DOCKSTRING pose archives, an external smina binary
and the Python 3.9 ODDT environment with its training tables. Its released scores and metrics are
frozen and checksum-recorded. The \texttt{refresh-public-analyses} target reruns a public subset;
it is not a universal producer pipeline.

The package starts from published score matrices. It does not rerun receptor preparation, search-box
selection, pose generation or docking, and those stages belong to the cited Docking-44 and
DOCKSTRING releases \cite{nikitin2025,dockstring2022}.

\sisubsection{S7.2 The reusable target-map audit}

\path{analysis/target_map_audit.py} applies the target-map checks of this paper to any new dense
score matrix. A test asserts that it imports nothing beyond the standard library, NumPy and
pandas, and that its source contains no reference to this project's evidence builder or to any
named dataset, so it carries no dependency on the manuscript.

The input is one table with a \texttt{.csv}, \texttt{.tsv}, \texttt{.csv.gz} or \texttt{.tsv.gz}
suffix; the separator follows from the suffix and any other suffix is rejected. The table needs
at least 4 rows, no duplicate column names and at least 3 target columns. Selection of the target
columns is required and takes exactly one of two forms, an explicit ordered list or a regular
expression searched against the column names. Table~\ref{tab:s7cli} is the complete command-line
interface.

\begin{table}[htbp]
\centering
\caption{Every option \texttt{analysis/target\_map\_audit.py} implements. The parser accepts no
others.}
\label{tab:s7cli}
\small
\begin{tabular}{llp{6.6cm}}
\toprule
Option & Default & Effect \\
\midrule
\texttt{input} (positional) & required & dense CSV/TSV table, optionally gzipped \\
\texttt{-{}-output-dir} & required & destination; must be absent or empty \\
\texttt{-{}-target-columns} & one of two required & ordered comma-separated target names \\
\texttt{-{}-target-regex} & one of two required & pattern searched against column names \\
\texttt{-{}-ligand-id} & none & identifier column; must have no missing value \\
\texttt{-{}-domain-column} & none & complete numeric column with non-zero range \\
\texttt{-{}-imputation} & \texttt{error} & \texttt{error}, \texttt{target-mean} or
\texttt{target-median} \\
\texttt{-{}-domain-quantiles} & none & number of quantile bins; requires a domain column, and a
domain column requires it \\
\texttt{-{}-pilot-sizes} & none & comma-separated row counts for held-out recovery \\
\texttt{-{}-pilot-repeats} & 100 & replicates per pilot size \\
\texttt{-{}-top-edge-count} & 10 & how many strongest positive and negative edges are compared \\
\texttt{-{}-panel-k} & none & number of medoid targets under distance $1-|r|$ \\
\texttt{-{}-seed} & 0 & seed for the pilot row permutations \\
\bottomrule
\end{tabular}
\end{table}

The tool fails closed. A missing target cell stops the run unless an imputation policy is chosen,
a nonnumeric cell always stops it and the message names up to 5 offending cells, and a
zero-variance target column stops it. All numerical work and validation finish before the output
directory is created, so a violated contract leaves no partial audit on disk. Errors print to
standard error with the prefix \texttt{target-map-audit: error:} and exit with status 2.

Six files are always written. \path{preprocessing_contract.json} records the schema version, the
SHA-256 of the script that produced the run, the input path with its byte count and digest, the
selected columns, the missing-data policy with any imputed values, the seed, the definition of
each surface, and a numerical check that row centring leaves every within-row pairwise target
difference unchanged. \path{spectra.json} holds, for the raw and the two-way-centred surface,
the participation-ratio dimension together with its exact mean-squared-correlation identity and
the absolute discrepancy between the two, the entropy effective dimension, the PC1 variance
fraction, the number of components reaching 90\% of trace, the mean and mean absolute edge
correlation, the negative-edge fraction, and the full eigenspectrum.
\path{target_correlation_raw.csv} and \path{target_correlation_two_way_centered.csv} hold the
square maps, \path{signed_edges.csv} the upper-triangle edge table under both transforms, and
\path{output_checksums.json} the byte count and SHA-256 of every file written.

Three optional analyses add files only when requested. \texttt{-{}-domain-quantiles} writes
\path{domain_bin_edges.csv} and \path{domain_map_comparisons.csv}, comparing every pair of
quantile bins under both transforms; bin boundaries come from NumPy quantiles, a value equal to
an interior boundary joins the higher bin, and each bin needs at least 3 rows.
\texttt{-{}-pilot-sizes} writes \path{pilot_recovery_replicates.csv} and
\path{pilot_recovery_summary.csv}: each replicate permutes the rows, splits them into a pilot
block and its row-disjoint complement, and compares the two maps on edge-rank Spearman, sign
agreement, mean absolute and root-mean-squared edge difference, and top positive and top negative
edge precision. Each replicate stores a SHA-256 of its sampled row indices, and the summary
reports the mean, the median and the 2.5--97.5\% range. \texttt{-{}-panel-k} writes
\path{target_panel_medoids.csv} and \path{target_panel_summary.json} from a deterministic
farthest-first initialisation followed by PAM swaps on distance $1-|r|$; both files carry the
warning that the result is not a biologically validated target panel.

A worked invocation on the shipped 20-ligand synthetic fixture, which has 4 target columns and a
molecular-weight column, is:

\begin{verbatim}
.venv/bin/python analysis/target_map_audit.py \
  analysis/fixtures/target_map_audit_synthetic.csv \
  --output-dir results/my_target_map_audit \
  --target-regex '^target_' --ligand-id ligand_id \
  --domain-column molecular_weight --domain-quantiles 2 \
  --pilot-sizes 8 --pilot-repeats 20 \
  --top-edge-count 3 --panel-k 2 --seed 20260810
\end{verbatim}

It writes 12 files and prints a JSON summary reporting 20 rows, 4 targets, a raw
participation-ratio dimension of 1.757 and a two-way-centred dimension of 1.850. The fixture is
built so that the residual map reverses between the light and heavy halves of the library: the
test suite asserts an edge-rank correlation below $-0.8$ and sign agreement below 0.5 between the
two molecular-weight bins. That is the support-dependence effect of the main text written into a
synthetic invariant.

The tool ships with \path{analysis/test_target_map_audit.py}, 7 tests that pass in about 1
second. They check the import restriction, the exactness of the participation-ratio identity
together with rank invariance under row centring, a complete run whose recorded checksums match
the files on disk, byte-for-byte determinism across two runs at the same seed, explicit column
selection from a gzipped TSV, fail-closed behaviour on missing and malformed cells, and rejection
of an invalid role column, an incomplete optional-analysis specification and a non-empty output
directory. The last of these also asserts that a pre-existing user file in that directory is left
untouched. \path{analysis/TARGET_MAP_AUDIT.md} is the user-facing documentation.
